%% ============================================================
%%  Chapter 6 — Release 3: Advanced Services and AI Features
%% ============================================================
\begingroup
\renewcommand{\baselinestretch}{0.85}
\chapter{Release 3: Advanced Services and AI Features}
\minitoc
\endgroup
\newpage

%% ------------------------------------------------------------
\section*{Introduction}
\addcontentsline{toc}{section}{Introduction}
%% ------------------------------------------------------------

This chapter presents the implementation of the third and final release of Tesla Academy.
Release 3 corresponds to Sprints 5 and 6 and introduces the advanced services that
differentiate the platform from conventional LMS solutions: online payment processing,
live session management, and a comprehensive suite of AI-powered educational features.

Sprint 5 covers paid course purchases, payment verification, payment history, and live
session scheduling, start, and participation. Sprint 6 introduces the contextual AI
learning assistant, AI-assisted assignment feedback, automatic question generation,
AI-powered visual theme generation, global platform indicators for the Super Administrator,
and the final platform validation and deployment.

%% ------------------------------------------------------------
\section{Sprint Backlog}
%% ------------------------------------------------------------

\begin{table}[H]
\centering
\caption{Sprint Backlog — Release 3}
\label{tab:backlog_r3}
\renewcommand{\arraystretch}{1.25}
\footnotesize
\begin{tabularx}{\textwidth}{|p{1.4cm}|p{1.5cm}|p{2.8cm}|p{6.8cm}|p{1.5cm}|}
\hline
\textbf{Sprint} & \textbf{Code} & \textbf{Feature} & \textbf{User Story} & \textbf{Est. (h)} \\
\hline
\multirow{7}{*}{Sprint 5}
  & US-042 & Buy course & Purchase a paid course from the platform & 24 \\
  \cline{2-5}
  & US-043 & Payment confirmation & Confirm payment and activate course enrollment & 16 \\
  \cline{2-5}
  & US-044 & Payment history & Consult payment history & 8 \\
  \cline{2-5}
  & US-045 & Schedule live session & Schedule a live session with date, time, and duration & 16 \\
  \cline{2-5}
  & US-046 & Start live session & Start a live session using the live session service & 16 \\
  \cline{2-5}
  & US-047 & Join live session & Join a live session from the student workspace & 8 \\
  \cline{2-5}
  & US-048 & Automatic session summary & Generate a summary after a live session & 16 \\
\hline
\multirow{7}{*}{Sprint 6}
  & US-049 & Contextual AI assistant & Ask the AI assistant questions related to course content & 24 \\
  \cline{2-5}
  & US-050 & Assignment feedback with AI & Receive AI-assisted feedback on assignments & 16 \\
  \cline{2-5}
  & US-051 & Question generation & Generate quiz questions automatically & 16 \\
  \cline{2-5}
  & US-052 & AI theme generation & Generate a visual theme from a text description & 16 \\
  \cline{2-5}
  & US-053 & Theme validation & Validate the generated theme before applying it & 8 \\
  \cline{2-5}
  & US-054 & Global indicators & Consult global platform indicators & 16 \\
  \cline{2-5}
  & US-055 & Final validation and deployment & Perform final validation and deploy the platform & 16 \\
\hline
\end{tabularx}
\end{table}

%% ------------------------------------------------------------
\section{Functional and Technical Specifications}
%% ------------------------------------------------------------

\subsection{Payment Processing (US-042, US-043, US-044)}

\begin{table}[H]
\centering
\caption{Specification — Paid Course Purchase}
\label{tab:payment_spec}
\renewcommand{\arraystretch}{1.2}
\begin{tabularx}{\textwidth}{|p{3.5cm}|X|}
\hline
\textbf{User Stories} & US-042, US-043, US-044 — Course purchase, payment confirmation, payment history. \\
\hline
\textbf{Business Rules} &
\begin{itemize}[noitemsep, topsep=2pt]
  \item Paid course enrollment requires a successful payment transaction;
  \item Two payment providers are supported: \textbf{Konnect} and \textbf{Flouci},
        both supporting Tunisian Dinar (TND);
  \item Each organization configures its own payment credentials in its
        \texttt{OrganizationConfig};
  \item Payment amounts are expressed in millimes (1 TND = 1000 millimes);
  \item An \texttt{Order} record is created upon checkout initiation with status
        \texttt{PENDING}, and updated to \texttt{SUCCESS} or \texttt{FAILED} after
        payment gateway callback verification;
  \item Upon confirmed payment, the student is automatically enrolled in the purchased
        course via a \texttt{BatchEnrollment} record.
\end{itemize} \\
\hline
\textbf{Technical Specification} &
\begin{itemize}[noitemsep, topsep=2pt]
  \item Konnect integration uses the \texttt{/payments/init-payment} API endpoint to
        generate a payment URL; verification is performed via the
        \texttt{/payments/:ref} endpoint;
  \item Flouci integration uses a public/private key pair stored per organization;
  \item The payment mode is determined by the \texttt{paymentMode} field in
        \texttt{OrganizationConfig}: \texttt{per\_course} or \texttt{subscription};
  \item Callback URLs include the \texttt{orderId} parameter, allowing the backend to
        match and update the corresponding \texttt{Order} record;
  \item Enrollment is granted by the \texttt{grantOrderEntitlements} function only after
        payment status is confirmed as \texttt{SUCCESS}.
\end{itemize} \\
\hline
\textbf{Acceptance Criteria} &
\begin{itemize}[noitemsep, topsep=2pt]
  \item A student can initiate checkout for a paid course;
  \item After successful payment, the student is enrolled and can access course content;
  \item A failed or cancelled payment does not grant enrollment;
  \item The student can view their payment history from their dashboard.
\end{itemize} \\
\hline
\end{tabularx}
\end{table}

\subsection{Live Session Management (US-045 to US-048)}

\begin{table}[H]
\centering
\caption{Specification — Live Sessions}
\label{tab:livekit_spec}
\renewcommand{\arraystretch}{1.2}
\begin{tabularx}{\textwidth}{|p{3.5cm}|X|}
\hline
\textbf{User Stories} & US-045 to US-048 — Schedule, start, join, and summarize live sessions. \\
\hline
\textbf{Business Rules} &
\begin{itemize}[noitemsep, topsep=2pt]
  \item An instructor schedules a session by providing a title, course association,
        date, time, and expected duration;
  \item Enrolled students receive an email notification when a new session is scheduled;
  \item The instructor starts the session from their workspace, which transitions the
        \texttt{Schedule} status to \texttt{LIVE} and issues a host-role LiveKit token;
  \item Students join using a participant-role token; they can view and hear but cannot
        publish audio/video by default unless the host grants permission;
  \item After the session ends, if transcription data is available, the AI microservice
        automatically generates a structured summary of the session.
\end{itemize} \\
\hline
\textbf{Technical Specification} &
\begin{itemize}[noitemsep, topsep=2pt]
  \item Live sessions use \textbf{LiveKit}, an open-source WebRTC infrastructure;
  \item A \texttt{LiveKit JWT token} is generated server-side using the
        \texttt{LIVEKIT\_API\_KEY} and \texttt{LIVEKIT\_API\_SECRET} environment
        variables;
  \item Host tokens carry \texttt{canPublish: true} and \texttt{canSubscribe: true};
        student tokens carry \texttt{canPublish: false};
  \item Session transcription is collected incrementally on the client and sent to the
        backend after the session ends via the \texttt{transcriptDraft} field;
  \item The AI summarization endpoint (\texttt{POST /internal/v1/schedules/summarize})
        receives the transcript and returns a structured summary stored in the
        \texttt{aiSummary} field of the \texttt{Schedule} record;
  \item A collaborative whiteboard is available per session via the
        \texttt{ScheduleWhiteboard} entity.
\end{itemize} \\
\hline
\textbf{Acceptance Criteria} &
\begin{itemize}[noitemsep, topsep=2pt]
  \item An instructor can schedule and start a session; students can join using the
        provided room token;
  \item Session participants can communicate in real time;
  \item After the session, a summary is automatically generated and displayed;
  \item Completed sessions cannot be restarted.
\end{itemize} \\
\hline
\end{tabularx}
\end{table}

\subsection{Contextual AI Assistant (US-049)}

\begin{table}[H]
\centering
\caption{Specification — Contextual AI Assistant (RAG Chat)}
\label{tab:rag_spec}
\renewcommand{\arraystretch}{1.2}
\begin{tabularx}{\textwidth}{|p{3.5cm}|X|}
\hline
\textbf{User Story} & US-049 — As a student, I want to ask the AI assistant questions about course content so that I can get contextual explanations. \\
\hline
\textbf{Business Rules} &
\begin{itemize}[noitemsep, topsep=2pt]
  \item The AI assistant is scoped to the content of the course the student is
        currently viewing;
  \item Responses are grounded in the indexed course content through RAG;
  \item Each organization uses its own OpenRouter API key for AI calls;
  \item Conversation history is maintained within the session for multi-turn coherence.
\end{itemize} \\
\hline
\textbf{Technical Specification} &
\begin{itemize}[noitemsep, topsep=2pt]
  \item When educational content (YouTube video, PDF, text) is published, the backend
        triggers the \texttt{POST /internal/v1/content/index} endpoint to extract text
        and store vector embeddings in the \texttt{content\_embeddings} pgvector table;
  \item When the student sends a message, the AI microservice computes a query
        embedding, retrieves the top-$k$ most similar content chunks via cosine
        similarity, and injects them into the system prompt;
  \item The language model call is made through OpenRouter using the configured model
        (default: \texttt{google/gemini-flash-1.5});
  \item The response includes metadata: \texttt{ragChunksUsed} and \texttt{ragStatus},
        allowing the frontend to indicate whether retrieval was successful.
\end{itemize} \\
\hline
\textbf{Acceptance Criteria} &
\begin{itemize}[noitemsep, topsep=2pt]
  \item A student can ask a question and receive a contextually relevant answer;
  \item The assistant references information from the indexed course content;
  \item The conversation history is maintained within the session.
\end{itemize} \\
\hline
\end{tabularx}
\end{table}

\subsection{AI-Assisted Assignment Feedback (US-050)}

\begin{table}[H]
\centering
\caption{Specification — AI Assignment Feedback}
\label{tab:ai_feedback_spec}
\renewcommand{\arraystretch}{1.2}
\begin{tabularx}{\textwidth}{|p{3.5cm}|X|}
\hline
\textbf{User Story} & US-050 — As a student, I want to receive automatic AI-assisted feedback on my assignment so that I can understand my mistakes. \\
\hline
\textbf{Business Rules} &
\begin{itemize}[noitemsep, topsep=2pt]
  \item AI feedback is generated automatically when a student submits an assignment;
  \item The feedback analyzes each question response against the correct answer and
        the defined evaluation criteria;
  \item AI feedback is stored alongside the submission and displayed to the student;
  \item AI feedback does not replace the instructor's grade; it is supplementary.
\end{itemize} \\
\hline
\textbf{Technical Specification} &
\begin{itemize}[noitemsep, topsep=2pt]
  \item Upon submission, the backend compiles a structured prompt containing the
        assignment questions, the student's answers, and the correct answers;
  \item The prompt is sent to \texttt{POST /internal/v1/assignments/feedback};
  \item The AI response is expected as structured JSON containing per-question feedback
        and an overall assessment;
  \item The result is stored in the \texttt{AssignmentAiFeedback} table and linked
        to the submission;
  \item Weak concept flags are extracted from the feedback and stored in
        \texttt{WeakConceptFlag} for long-term student progress analysis.
\end{itemize} \\
\hline
\textbf{Acceptance Criteria} &
\begin{itemize}[noitemsep, topsep=2pt]
  \item After submitting an assignment, the student sees an AI-generated feedback
        section alongside their submission;
  \item The feedback addresses each question individually;
  \item The feedback is generated asynchronously and does not block the submission flow.
\end{itemize} \\
\hline
\end{tabularx}
\end{table}

\subsection{AI Question Generation (US-051)}

\begin{table}[H]
\centering
\caption{Specification — AI Question Generation}
\label{tab:qgen_spec}
\renewcommand{\arraystretch}{1.2}
\begin{tabularx}{\textwidth}{|p{3.5cm}|X|}
\hline
\textbf{User Story} & US-051 — As an instructor, I want to generate quiz questions automatically so that I can save time when preparing assessments. \\
\hline
\textbf{Business Rules} &
\begin{itemize}[noitemsep, topsep=2pt]
  \item The instructor specifies a topic, difficulty level, number of questions, and
        question type;
  \item The platform generates the questions using the LLM and presents them for review
        before insertion into the assignment;
  \item The instructor can edit, discard, or accept individual generated questions.
\end{itemize} \\
\hline
\textbf{Technical Specification} &
\begin{itemize}[noitemsep, topsep=2pt]
  \item A structured prompt is sent to
        \texttt{POST /internal/v1/assignments/generate-questions};
  \item The response is expected as JSON matching the internal
        \texttt{AssignmentQuestion} schema;
  \item A validation layer (\texttt{validateGeneratedAssignmentQuestion}) ensures that
        generated questions conform to the schema before they are offered to the instructor.
\end{itemize} \\
\hline
\textbf{Acceptance Criteria} &
\begin{itemize}[noitemsep, topsep=2pt]
  \item The instructor can request question generation by filling a short form;
  \item The generated questions appear in the assignment editor for review;
  \item The instructor can selectively add generated questions to the assignment.
\end{itemize} \\
\hline
\end{tabularx}
\end{table}

\subsection{AI Visual Theme Generation (US-052, US-053)}

\begin{table}[H]
\centering
\caption{Specification — AI Theme Generation}
\label{tab:theme_spec}
\renewcommand{\arraystretch}{1.2}
\begin{tabularx}{\textwidth}{|p{3.5cm}|X|}
\hline
\textbf{User Stories} & US-052, US-053 — Generate and validate an AI-powered visual theme. \\
\hline
\textbf{Business Rules} &
\begin{itemize}[noitemsep, topsep=2pt]
  \item The Organization Administrator provides a textual brand description as input;
  \item The AI microservice generates a complete CSS variable theme for the organization;
  \item The administrator previews the generated theme before applying it;
  \item Applying the theme updates the \texttt{customCss} field in \texttt{OrganizationConfig},
        which is injected dynamically into the student-facing frontend.
\end{itemize} \\
\hline
\textbf{Technical Specification} &
\begin{itemize}[noitemsep, topsep=2pt]
  \item The \texttt{generate\_theme\_bundle} function in \texttt{theme\_engine.py}
        constructs a prompt and calls OpenRouter to obtain a JSON payload containing
        \texttt{primaryColor}, \texttt{secondaryColor}, \texttt{fontFamily}, and
        \texttt{customCss};
  \item Generated CSS is expressed entirely in \texttt{oklch()} color format to ensure
        consistency with the Tailwind CSS v4 design token system;
  \item A \texttt{\_sanitize\_css} function validates the output against forbidden
        patterns (external URLs, \texttt{@import}, \texttt{expression()}), then merges
        the AI output with a deterministic scaffold to guarantee all required CSS
        variables are present;
  \item The final theme includes both \texttt{:root} (light mode) and \texttt{.dark}
        (dark mode) variable blocks;
  \item All 51 required CSS variables defined in \texttt{REQUIRED\_THEME\_VARS} must
        be present in the sanitized output; missing variables are filled from the
        scaffold.
\end{itemize} \\
\hline
\textbf{Acceptance Criteria} &
\begin{itemize}[noitemsep, topsep=2pt]
  \item The administrator enters a brand description and receives a generated theme;
  \item A live preview shows the theme applied to a sample of the platform UI;
  \item Clicking apply updates the organization's visual appearance across the platform;
  \item Invalid or unsafe CSS is never applied.
\end{itemize} \\
\hline
\end{tabularx}
\end{table}

%% ------------------------------------------------------------
\section{Design}
%% ------------------------------------------------------------

\subsection{Static Modeling}

\subsubsection{Class Diagram}

Figure~\ref{fig:class_r3} presents the class diagram for Release 3, focusing on the
payment, live session, and AI feature entities.

\begin{figure}[H]
  \centerline{\includegraphics[width=\textwidth]{class_diagram_r3.png}}
  \caption{Class Diagram — Release 3}
  \label{fig:class_r3}
\end{figure}

The main entities and relationships introduced in this release are:

\begin{itemize}
  \item An \texttt{Order} links a \texttt{User} to a purchased entity (a \texttt{Batch}
        or \texttt{TestSeries}) and tracks the payment status and provider;
  \item A \texttt{Schedule} (live session) belongs to an \texttt{Organization} and is
        optionally associated with a \texttt{Batch}; it tracks the LiveKit room name,
        session status, and transcript;
  \item A \texttt{ScheduleWhiteboard} is a collaborative drawing board associated with
        one \texttt{Schedule};
  \item An \texttt{AssignmentAiFeedback} record stores the AI-generated feedback for
        one \texttt{AssignmentSubmission};
  \item A \texttt{WeakConceptFlag} captures topics identified as weak for a student,
        derived from AI feedback analysis;
  \item The \texttt{OrganizationConfig} entity stores the \texttt{customCss} field
        used to inject the AI-generated theme, as well as payment credentials and
        the OpenRouter API key.
\end{itemize}

\subsection{Dynamic Modeling}

\subsubsection{Payment Flow Sequence Diagram}

Figure~\ref{fig:seq_payment} illustrates the paid course purchase flow. The student
initiates checkout from the course detail page. The backend creates an \texttt{Order}
with status \texttt{PENDING} and calls the payment gateway API to obtain a payment URL.
The student is redirected to the gateway's hosted checkout page. After payment, the
gateway redirects the student to the platform's callback URL. The backend verifies the
payment status via the gateway API and, on success, transitions the order to
\texttt{SUCCESS} and grants the course enrollment.

\begin{figure}[H]
  \centerline{\includegraphics[width=0.95\textwidth]{seq_payment.png}}
  \caption{Paid Course Purchase Sequence Diagram}
  \label{fig:seq_payment}
\end{figure}

\subsubsection{Live Session Sequence Diagram}

Figure~\ref{fig:seq_live} illustrates the live session workflow. The instructor schedules
a session, which triggers enrollment email notifications. At session time, the instructor
clicks \textit{Start}. The backend generates a host-role LiveKit token and updates the
session status to \texttt{LIVE}. Students request a participant token and connect to the
same LiveKit room. After the session ends, the transcript is submitted and the AI
microservice generates a summary.

\begin{figure}[H]
  \centerline{\includegraphics[width=0.95\textwidth]{seq_live_session.png}}
  \caption{Live Session Sequence Diagram}
  \label{fig:seq_live}
\end{figure}

\subsubsection{RAG Chat Sequence Diagram}

Figure~\ref{fig:seq_rag} illustrates the Retrieval-Augmented Generation (RAG) pipeline
used by the contextual AI assistant. When a student sends a message, the AI microservice
generates a vector embedding of the query, searches the \texttt{content\_embeddings}
table for the most semantically similar content chunks, injects these chunks into the
system prompt, and calls the language model via OpenRouter to produce a grounded response.

\begin{figure}[H]
  \centerline{\includegraphics[width=0.9\textwidth]{seq_rag_chat.png}}
  \caption{RAG Chat Sequence Diagram}
  \label{fig:seq_rag}
\end{figure}

\subsubsection{AI Theme Generation Sequence Diagram}

Figure~\ref{fig:seq_theme} illustrates the AI theme generation flow. The administrator
submits a brand description through the theme editor. The backend forwards the request to
the AI microservice, which constructs a structured prompt, calls OpenRouter, parses and
sanitizes the CSS output, and returns the theme bundle. The administrator previews the
result and optionally applies it, which stores the \texttt{customCss} in the organization
configuration and applies it immediately to the student-facing interface.

\begin{figure}[H]
  \centerline{\includegraphics[width=0.9\textwidth]{seq_theme_generation.png}}
  \caption{AI Theme Generation Sequence Diagram}
  \label{fig:seq_theme}
\end{figure}

%% ------------------------------------------------------------
\section{Implementation}
%% ------------------------------------------------------------

\subsection{Payment Integration}

Tesla Academy integrates two Tunisian payment gateways: \textbf{Konnect} and
\textbf{Flouci}. Both gateways operate with TND-denominated transactions expressed in
millimes. The active gateway for each organization is determined by the
\texttt{paymentMode} field in its configuration.

The checkout flow creates a pending \texttt{Order} record, calls the selected gateway
API to generate a hosted payment URL, and redirects the student. The callback handler
verifies the payment reference with the gateway and, on confirmed success, calls
\texttt{grantOrderEntitlements} to create the enrollment. The student is then redirected
to the newly accessible course.

\begin{figure}[H]
  \centerline{\includegraphics[width=\textwidth]{screenshot_checkout.png}}
  \caption{Course Checkout Page — Student View}
  \label{fig:sc_checkout}
\end{figure}

\begin{figure}[H]
  \centerline{\includegraphics[width=\textwidth]{screenshot_payment_history.png}}
  \caption{Payment History — Student Dashboard}
  \label{fig:sc_payment_history}
\end{figure}

\subsection{Live Session Room}

The live session room is built on top of the LiveKit SDK. The instructor's view provides
controls for microphone, camera, screen sharing, and participant management. The student
view provides a viewer layout with audio/video reception and a real-time chat panel. A
collaborative whiteboard panel is also accessible during the session.

Transcription is collected progressively on the client by concatenating speech recognition
segments. After the host ends the session, the accumulated transcript is submitted to the
backend and forwarded to the AI microservice for summary generation. The resulting summary
is stored and displayed to both the instructor and enrolled students.

\begin{figure}[H]
  \centerline{\includegraphics[width=\textwidth]{screenshot_live_session.png}}
  \caption{Live Session Room — Instructor View}
  \label{fig:sc_live_session}
\end{figure}

\begin{figure}[H]
  \centerline{\includegraphics[width=\textwidth]{screenshot_session_summary.png}}
  \caption{AI-Generated Session Summary}
  \label{fig:sc_session_summary}
\end{figure}

\subsection{Contextual AI Learning Assistant}

The AI assistant is accessible to students from within any course page as a side panel.
The student types a question in natural language; the assistant retrieves relevant content
chunks from the course materials using vector similarity search and constructs a
contextually grounded response.

The content indexing pipeline is triggered automatically when an instructor publishes
a content item. For YouTube videos, the transcript is extracted using the YouTube
Transcript API. For PDF documents, text is extracted with \texttt{pypdf}. The extracted
text is chunked and stored as vector embeddings in the \texttt{content\_embeddings} table
via the \texttt{pgvector} extension.

\begin{figure}[H]
  \centerline{\includegraphics[width=\textwidth]{screenshot_ai_assistant.png}}
  \caption{Contextual AI Learning Assistant — Student View}
  \label{fig:sc_ai_assistant}
\end{figure}

\subsection{AI-Assisted Assignment Feedback and Question Generation}

When a student submits an assignment, the backend triggers an asynchronous AI feedback
request. The feedback analyzes each response against the expected answer and the assignment
grading criteria. The result is displayed to the student in a structured panel alongside
their submission.

The question generation interface allows instructors to specify a topic, question type,
difficulty, and desired count. The AI returns a set of structured questions that the
instructor can review, edit, and add to an assignment with a single click.

\begin{figure}[H]
  \centerline{\includegraphics[width=\textwidth]{screenshot_ai_feedback.png}}
  \caption{AI Assignment Feedback Panel — Student View}
  \label{fig:sc_ai_feedback}
\end{figure}

\begin{figure}[H]
  \centerline{\includegraphics[width=\textwidth]{screenshot_question_gen.png}}
  \caption{AI Question Generation — Instructor View}
  \label{fig:sc_question_gen}
\end{figure}

\subsection{AI Visual Theme Generation}

The theme generator allows the Organization Administrator to describe their brand
identity in free text. The AI microservice constructs a deterministic CSS token scaffold
based on the extracted color keywords and font preferences, then asks the language model
to refine it. The resulting theme defines all 51 required CSS variables in
\texttt{oklch()} color space for both light and dark modes.

The generated CSS is validated to reject any forbidden patterns (external URLs,
executable JavaScript, Tailwind at-rules), merged with the scaffold to ensure
completeness, and returned to the frontend for preview. Upon administrator approval, the
theme is saved in the organization configuration and injected into the platform via a
\texttt{<style>} tag in the HTML head.

\begin{figure}[H]
  \centerline{\includegraphics[width=\textwidth]{screenshot_theme_editor.png}}
  \caption{AI Theme Generation — Organization Administrator View}
  \label{fig:sc_theme_editor}
\end{figure}

\subsection{Super Administrator Global Dashboard}

The Super Administrator dashboard provides a global view of all organizations registered
on the platform. Key indicators displayed include the total number of active organizations,
total registered users across all tenants, total published courses, and platform-wide
revenue. A per-organization breakdown allows the Super Administrator to monitor individual
organization activity and manage activation status.

\begin{figure}[H]
  \centerline{\includegraphics[width=\textwidth]{screenshot_superadmin_dashboard.png}}
  \caption{Global Indicators — Super Administrator Dashboard}
  \label{fig:sc_superadmin_dashboard}
\end{figure}

%% ------------------------------------------------------------
\section*{Conclusion}
\addcontentsline{toc}{section}{Conclusion}
%% ------------------------------------------------------------

This chapter presented the implementation of Release 3 and completed the full description
of the Tesla Academy platform. The third release introduced the advanced services that
define the platform's competitive positioning: Tunisian payment gateway integration with
Konnect and Flouci, real-time live sessions powered by LiveKit with automatic AI
summarization, and a comprehensive AI feature suite including a RAG-based contextual
learning assistant, AI-assisted assignment feedback, automatic question generation, and
AI-powered visual theme generation.

Together, the three releases of Tesla Academy deliver a complete, multi-tenant, and
AI-enhanced Learning Management System adapted to the needs of Tunisian and regional
educational institutions. The platform is designed to be scalable, maintainable, and
extensible, providing a solid foundation for continued development and market deployment.